\chapter{Opérateurs compacts et théorème spectral}
\label{ch:b3:spectral}

La diagonalisation est le joyau de l'algèbre linéaire de
dimension finie~: une matrice symétrique possède une \omterm{def:b3:topology:basis}{base}
orthonormée de vecteurs propres. En dimension infinie ceci
échoue pour les \omterm{def:b3:spectral:selfadjoint}{opérateurs autoadjoints} bornés en général ---
la multiplication par $x$ sur $L^2(\intcc01)$ n'a
\emph{aucune} valeur propre (\cref{exo:b3:spectral:6}) ---
mais survit, sous une forme parfaite, pour les opérateurs
\emph{presque de dimension finie}~: les \omterm{def:b3:spectral:compact}{compacts}. Le théorème
spectral pour les \omterm{def:b3:spectral:compact}{opérateurs compacts} \omterm{def:b3:spectral:selfadjoint}{autoadjoints} est le
théorème le plus utilisé de l'analyse fonctionnelle
appliquée~: il diagonalise les équations intégrales, anime
l'\omterm{thm:b3:spectral:fredholm}{alternative de Fredholm}, et (problème du week-end) résout
la corde vibrante, produisant la \omterm{def:b3:topology:basis}{base} en sinus de l'analyse
de Fourier à partir de la pure théorie des opérateurs ---
avec $\zeta(2) = \frac{\pi^2}6$ d'Euler tombant d'une
formule de trace en cadeau de départ. Tout au long, $H$ est
un \omterm{def:b3:hilbert:inner}{espace de Hilbert} sur $\C$ (ou $\R$~; les énoncés
s'adaptent), et les opérateurs sont bornés.

\section{Opérateurs compacts}

\begin{definition}\label{def:b3:spectral:compact}
$T \in \mathcal L(E, F)$ ($E, F$ de Banach) est
\emph{compact}\index{opérateur compact} si l'image $T(B)$ de
la boule unité est relativement compacte dans $F$ --- de
façon équivalente, toute suite bornée $(x_n)$ admet une
sous-suite avec $(Tx_{n_k})$ convergente. Les opérateurs de
rang fini sont compacts (ensembles bornés en dimension
finie)~; l'identité d'un espace de dimension infinie ne
l'est jamais (théorème de Riesz, deuxième année).
\end{definition}

\begin{proposition}\label{prop:b3:spectral:ideal}
Les \omterm{def:b3:spectral:compact}{opérateurs compacts} $\mathcal K(E, F)$ forment un
sous-espace fermé de $\mathcal L(E, F)$, et un \omterm{def:b3:rings:ideal}{idéal}
bilatère~: $S$ \omterm{def:b3:spectral:compact}{compact} $\Rightarrow$ $AS$ et $SB$ \omterm{def:b3:spectral:compact}{compacts}
pour $A, B$ bornés. De plus, dans un \omterm{def:b3:hilbert:inner}{espace de Hilbert}, tout
\omterm{def:b3:spectral:compact}{opérateur compact} est limite en \omterm{def:b3:banach:operator}{norme d'opérateurs} de rang
fini.
\end{proposition}

\begin{proof}
Sous-espace~: clair par la caractérisation séquentielle.
\omterm{def:b3:rings:ideal}{Idéal}~: les applications bornées envoient suites
convergentes sur suites convergentes et bornées sur bornées.
Fermeture~: soit $T_n \to T$ avec $T_n$ \omterm{def:b3:spectral:compact}{compact}, et $(x_k)$
bornée par $1$~; une extraction diagonale rend
$(T_nx_{k_j})_j$ convergente pour tout $n$~; alors
$(Tx_{k_j})$ est de Cauchy, car
\[
\norm{Tx_{k_j} - Tx_{k_l}} \leq 2\vertiii{T - T_n} +
\norm{T_nx_{k_j} - T_nx_{k_l}} ,
\]
en choisissant d'abord $n$ puis les indices. Approximation
dans les \omterm{def:b3:hilbert:inner}{espaces de Hilbert}~: soit $T$ \omterm{def:b3:spectral:compact}{compact}, $K =
\overline{T(B)}$ \omterm{def:b3:spectral:compact}{compact}~; pour $\varepsilon$ donné, recouvrir
$K$ par un nombre fini de boules $B(y_i, \varepsilon)$ et
soit $P$ la projection orthogonale sur $V =
\operatorname{Vect}(y_1, \dots, y_m)$ (fermé~: dimension
finie). Alors $PT$ est de rang fini, et pour $\norm x \leq
1$~: en prenant $y_i$ avec $\norm{Tx - y_i} < \varepsilon$,
\[
\norm{Tx - PTx} \leq \norm{Tx - y_i} + \norm{P(y_i - Tx)}
\leq 2\varepsilon
\]
($y_i = Py_i$~; $\vertiii P \leq 1$)~: $\vertiii{T - PT} \leq
2\varepsilon$.
\end{proof}

\begin{example}\label{ex:b3:spectral:examples}
(a) Opérateurs diagonaux sur $\ell^2$~: $T(x_n) =
(\lambda_nx_n)$ est \omterm{def:b3:spectral:compact}{compact} ssi $\lambda_n \to 0$
(\cref{exo:b3:spectral:2}).
(b) Opérateurs à noyau sur $\mathcal C(\intcc01)$~: \omterm{def:b3:spectral:compact}{compacts}
par Ascoli (\cref{exo:b3:complete:7}).
(c) \emph{Opérateurs de Hilbert--Schmidt}\index{opérateur de
Hilbert--Schmidt}~: pour $k \in L^2(\intcc01^2)$,
\[
(T_kf)(x) = \int_0^1k(x, y)\,f(y)\,\dd y
\]
définit un \omterm{def:b3:spectral:compact}{opérateur compact} sur $L^2(\intcc01)$ avec
$\vertiii{T_k} \leq \norm k_{L^2}$
(\cref{exo:b3:spectral:4}~: tronquer le développement en \omterm{def:b3:topology:basis}{base}
de $k$ exhibe $T_k$ comme limite d'opérateurs de rang fini).
\end{example}

\section{Opérateurs autoadjoints}

\begin{definition}\label{def:b3:spectral:selfadjoint}
$T \in \mathcal L(H)$ est \emph{autoadjoint}\index{opérateur
autoadjoint} si $T = T^*$ (\cref{exo:b3:hilbert:8}), i.e.\
$\langle Tx, y\rangle = \langle x, Ty\rangle$ pour tous $x,
y$. Alors $\langle x, Tx\rangle \in \R$ pour tout $x$ (égal
à son conjugué).
\end{definition}

\begin{proposition}\label{prop:b3:spectral:sanorm}
Pour $T$ \omterm{def:b3:spectral:selfadjoint}{autoadjoint}~\index{rayon numérique}~:
\[
\vertiii T = \sup_{\norm x \leq 1}\ \abs{\langle x,
Tx\rangle} .
\]
Les valeurs propres de $T$ sont réelles, et les vecteurs
propres pour des valeurs propres distinctes sont orthogonaux.
\end{proposition}

\begin{proof}
Soit $M$ le supremum~; $M \leq \vertiii T$ par
Cauchy--Schwarz. Réciproquement, l'identité de type
polarisation
\[
\langle x{+}y, T(x{+}y)\rangle - \langle x{-}y,
T(x{-}y)\rangle = 4\operatorname{Re}\langle y, Tx\rangle
\]
(développer~; les termes croisés $\langle y, Tx\rangle +
\langle x, Ty\rangle = 2\operatorname{Re}\langle y,
Tx\rangle$ par autoadjonction) donne, avec l'identité du
parallélogramme,
\[
4\operatorname{Re}\langle y, Tx\rangle \leq
M\bigl(\norm{x{+}y}^2 + \norm{x{-}y}^2\bigr) = 2M\bigl(\norm
x^2 + \norm y^2\bigr).
\]
Pour $\norm x = 1$ avec $Tx \neq 0$, prendre $y =
Tx/\norm{Tx}$~: $4\norm{Tx} \leq 4M$. Donc $\vertiii T \leq
M$. Valeurs propres~: $Tx = \lambda x$, $x \ne 0$ donne
$\lambda\norm x^2 = \langle x, Tx\rangle \in \R$.
Orthogonalité~: $\lambda\langle x, y\rangle = \langle Tx,
y\rangle = \langle x, Ty\rangle = \mu\langle x, y\rangle$
avec $\lambda \neq \mu$ réels.
\end{proof}

\section{Le théorème spectral}

\begin{lemma}[Existence d'une valeur propre extrémale]
\label{lem:b3:spectral:existence}
Soit $T \neq 0$ \omterm{def:b3:spectral:compact}{compact} et \omterm{def:b3:spectral:selfadjoint}{autoadjoint}. Alors $\vertiii T$
ou $-\vertiii T$ est une valeur propre de $T$.
\end{lemma}

\begin{proof}
Par la~\cref{prop:b3:spectral:sanorm}, prendre des vecteurs
unitaires $x_n$ avec $\langle x_n, Tx_n\rangle \to \mu$, où
$\abs\mu = \vertiii T > 0$ (passer à une sous-suite pour fixer
le signe). Alors
\[
\norm{Tx_n - \mu x_n}^2
= \norm{Tx_n}^2 - 2\mu\langle x_n, Tx_n\rangle + \mu^2
\leq \vertiii T^2 - 2\mu\langle x_n, Tx_n\rangle + \mu^2
\longrightarrow 2\mu^2 - 2\mu\cdot\mu = 0 .
\]
Par compacité, une sous-suite $Tx_{n_k} \to y$~; alors $\mu
x_{n_k} = Tx_{n_k} - (Tx_{n_k} - \mu x_{n_k}) \to y$, donc
$x_{n_k} \to x = y/\mu$, un vecteur unitaire, et la
continuité donne $Tx = \mu x$.
\end{proof}

\begin{theorem}[Théorème spectral pour les opérateurs
compacts autoadjoints]\label{thm:b3:spectral:spectral}
Soit $T$ un \omterm{def:b3:spectral:compact}{opérateur compact} \omterm{def:b3:spectral:selfadjoint}{autoadjoint} sur un \omterm{def:b3:hilbert:inner}{espace de
Hilbert} $H$.\index{théorème spectral (compact autoadjoint)}
\begin{enumerate}
  \item $H$ admet un système orthonormé $(e_n)_{n \in N}$
        ($N$ fini ou dénombrable) de vecteurs propres de $T$,
        avec valeurs propres réelles non nulles
        $(\lambda_n)$, tels que
        \[
        Tx = \sum_{n\in N}\lambda_n\,\langle e_n, x\rangle\,
        e_n \qquad (x \in H),
        \]
        et $H = \ker T \,\oplus^\perp\,
        \overline{\operatorname{Vect}}(e_n : n \in N)$.
  \item Si $N$ est infini, $\lambda_n \to 0$~; pour chaque
        $\delta > 0$ seuls un nombre fini de $n$ ont
        $\abs{\lambda_n} \geq \delta$, et chaque espace
        propre $\ker(T - \lambda)$, $\lambda \neq 0$, est de
        dimension finie.
  \item En complétant $(e_n)$ par une \omterm{def:b3:topology:basis}{base} orthonormée de
        $\ker T$ on obtient, lorsque $H$ est \omterm{def:b3:galois:separable}{séparable}, une
        \omterm{def:b3:topology:basis}{base} orthonormée de $H$ faite de vecteurs propres~:
        $T$ est diagonalisé.
\end{enumerate}
\end{theorem}

\begin{proof}
(2) d'abord. Si une infinité de vecteurs propres orthonormés
$x_k$ avaient $\abs{\lambda_{(k)}} \geq \delta$~:
$\norm{Tx_k - Tx_l}^2 = \lambda_{(k)}^2 + \lambda_{(l)}^2
\geq 2\delta^2$ (orthogonalité, Pythagore)~: aucune
sous-suite convergente de $(Tx_k)$, contredisant la
compacité de $T$ sur la suite bornée $(x_k)$. Ceci borne par
un nombre fini, pour chaque $\delta$, la multiplicités
totale des valeurs propres hors de $\intoo{-\delta}\delta$~;
la dénombrabilité et $\lambda_n \to 0$ s'ensuivent.

(1) Soit $H_0$ l'enveloppe fermée de \emph{tous} les vecteurs
propres à valeurs propres non nulles, organisés (par (2) et
Gram--Schmidt dans chaque espace propre de dimension finie,
orthogonalité entre espaces propres par
la~\cref{prop:b3:spectral:sanorm}) en un système orthonormé
$(e_n)$ de valeurs propres $\lambda_n \neq 0$. $T$ envoie
$H_0$ dans $H_0$, et aussi $H_0^\perp$ dans $H_0^\perp$~:
pour $y \perp H_0$ et $e$ un vecteur propre, $\langle e,
Ty\rangle = \langle Te, y\rangle = \lambda\langle e,
y\rangle = 0$. La restriction $T' = T\restriction_{H_0^\perp}$
est compacte autoadjointe sur l'\omterm{def:b3:hilbert:inner}{espace de Hilbert}
$H_0^\perp$~; si $T' \neq 0$,
\cref{lem:b3:spectral:existence} produit un vecteur propre
de $T$ à valeur propre non nulle dans $H_0^\perp$ ---
impossible, de tels vecteurs vivent dans $H_0$. Donc $T' =
0$~: $H_0^\perp \subseteq \ker T$. Réciproquement $\ker T
\perp$ chaque $e_n$ ($\langle e_n, z\rangle =
\frac1{\lambda_n}\langle Te_n, z\rangle =
\frac1{\lambda_n}\langle e_n, Tz\rangle = 0$)~: $\ker T
\subseteq H_0^\perp$, d'où $\ker T = H_0^\perp$ et la
\omterm{thm:b3:hilbert:decomposition}{décomposition orthogonale}. Le développement~: pour $x = z +
\sum_nc_ne_n$ ($z \in \ker T$, $c_n = \langle e_n,
x\rangle$~; \cref{thm:b3:hilbert:parseval}(1) sur $H_0$),
la continuité de $T$ donne $Tx = \sum_nc_n\lambda_ne_n$.

(3) $\ker T$, sous-espace fermé d'un espace \omterm{def:b3:galois:separable}{séparable}, est
\omterm{def:b3:galois:separable}{séparable}~: il a une \omterm{def:b3:topology:basis}{base} orthonormée
(\cref{prop:b3:hilbert:gramschmidt})~; l'union est une \omterm{def:b3:topology:basis}{base}
orthonormée de $H$ par la décomposition de (1).
\end{proof}

\begin{theorem}[Alternative de Fredholm]
\label{thm:b3:spectral:fredholm}
Soit $T$ \omterm{def:b3:spectral:compact}{compact} \omterm{def:b3:spectral:selfadjoint}{autoadjoint} et $\lambda \in \R\setminus
\{0\}$.\index{alternative de Fredholm}
\begin{enumerate}
  \item Si $\lambda$ n'est pas valeur propre, alors $T -
        \lambda I$ est bijectif d'inverse borné~: pour tout
        $f$, l'équation $Tx - \lambda x = f$ a exactement une
        solution, dépendant continûment de $f$.
  \item Si $\lambda$ est valeur propre, $Tx - \lambda x = f$
        est \omterm{def:b3:groups:derived}{résoluble} ssi $f \perp \ker(T - \lambda I)$, et
        la solution est unique à ce noyau (de dimension
        finie) près.
\end{enumerate}
\end{theorem}

\begin{proof}
Décomposer $x = z + \sum c_ne_n$ et $f = w + \sum d_ne_n$
selon le~\cref{thm:b3:spectral:spectral} ($z, w \in \ker T$).
L'équation se lit
\[
-\lambda z = w, \qquad (\lambda_n - \lambda)\,c_n = d_n\
(n \in N).
\]
(1) $\lambda \notin \{\lambda_n\}\cup\{0\}$~: par (2) du
théorème spectral, $\inf_n\abs{\lambda_n - \lambda} = \delta
> 0$ (les valeurs propres s'accumulent seulement en $0 \neq
\lambda$). Résoudre~: $z = -w/\lambda$, $c_n =
d_n/(\lambda_n - \lambda)$, avec $\sum\abs{c_n}^2 \leq
\delta^{-2}\sum\abs{d_n}^2$~: une solution unique avec $\norm
x \leq C\norm f$.
(2) $\lambda = \lambda_{n}$ pour $n$ dans un ensemble fini
$F$~: la résolubilité de $(\lambda_n - \lambda)c_n = d_n$
pour $n \in F$ exige $d_n = 0$, i.e.\ $f \perp e_n$ ($n \in
F$), i.e.\ $f \perp \ker(T - \lambda I)$~; les $c_n$, $n
\in F$, sont alors \omterm{def:b3:modules:free}{libres}.
\end{proof}

\begin{example}\label{ex:b3:spectral:minxy}
Sur $L^2(\intcc01)$, soit $Tf(x) = \int_0^1\min(x,
y)f(y)\dd y$~: un \omterm{ex:b3:spectral:examples}{opérateur de Hilbert--Schmidt} à noyau réel
symétrique~: \omterm{def:b3:spectral:compact}{compact} et \omterm{def:b3:spectral:selfadjoint}{autoadjoint}. Résoudre $Tf = \lambda
f$~: la relation $\bigl(Tf\bigr)(x) = \int_0^xyf(y)\dd y +
x\int_x^1f(y)\dd y$ montre que $u = Tf$ satisfait $u'' = -f$
(deux dérivations, légitimes pour $f$ \omterm{def:b3:topology:continuity}{continue}, et $Tf$ est
\omterm{def:b3:topology:continuity}{continue} pour $f \in L^2$~: convergence dominée), avec $u(0)
= 0$ et $u'(1) = 0$. Donc les fonctions propres résolvent
$\lambda u'' = -u$, $u(0) = 0$, $u'(1) = 0$~:
\[
u_n(x) = \sin\Bigl(\bigl(n + \tfrac12\bigr)\pi x\Bigr),
\qquad
\lambda_n = \frac{1}{\bigl(n + \frac12\bigr)^2\pi^2}
\quad (n \geq 0),
\]
et le théorème spectral affirme --- sans aucune théorie de
Fourier --- que ces sinus forment une \omterm{def:b3:topology:basis}{base} orthonormée de
$L^2(\intcc01)$ après normalisation (le noyau de $T$ est
$0$~: $Tf = 0$ force, par les deux dérivations, $f = 0$
p.p.). Le problème du week-end parcourt le même cercle
d'idées pour la corde vibrante et extrait $\zeta(2)$ de la
trace.
\end{example}

\begin{method}\label{met:b3:spectral:method}
Étant donnée une équation intégrale ou différentielle~: (1)
la reformuler comme $(I - \lambda K)u = f$ ou $Ku = \lambda
u$ avec $K$ un opérateur intégral~; (2) vérifier que $K$ est
\omterm{def:b3:spectral:compact}{compact} (noyau de Hilbert--Schmidt, ou Ascoli) et, si
possible, \omterm{def:b3:spectral:selfadjoint}{autoadjoint} (noyau réel symétrique)~; (3)
diagonaliser avec le théorème spectral ou invoquer
l'\omterm{thm:b3:spectral:fredholm}{alternative de Fredholm} pour la résolubilité~; (4) lire
l'existence, l'unicité, la stabilité, et les formules en
séries pour les solutions dans la \omterm{def:b3:topology:basis}{base} propre. Les opérateurs
différentiels sont non bornés, mais leurs \emph{inverses}
(opérateurs de Green) sont \omterm{def:b3:spectral:compact}{compacts}~: toujours inverser
d'abord.
\end{method}

\section{Exercices}

\begin{exercise}[$\star$]\label{exo:b3:spectral:1}
(a) Montrer qu'un opérateur borné à image de dimension finie
est \omterm{def:b3:spectral:compact}{compact}.
(b) Montrer que l'identité d'un espace normé est compacte ssi
la dimension est finie (Riesz, deuxième année). En déduire
qu'un \omterm{def:b3:spectral:compact}{opérateur compact} sur un espace de dimension infinie
n'est jamais inversible d'inverse borné.
\end{exercise}

\begin{exercise}[$\star$]\label{exo:b3:spectral:2}
Soit $T(x_1, x_2, \dots) = (\lambda_1x_1, \lambda_2x_2,
\dots)$ sur $\ell^2$, avec $(\lambda_n)$ bornée.
(a) Montrer $\vertiii T = \sup\abs{\lambda_n}$.
(b) Montrer que $T$ est \omterm{def:b3:spectral:compact}{compact} ssi $\lambda_n \to 0$.
\emph{(Pour $\Leftarrow$, tronquer~; pour $\Rightarrow$,
tester sur $(e_n)$.)}
(c) Quand $T$ est-il \omterm{def:b3:spectral:selfadjoint}{autoadjoint}~? Vérifier le théorème
spectral par inspection dans ce cas.
\end{exercise}

\begin{exercise}[$\star\star$]\label{exo:b3:spectral:3}
Donner les détails de la propriété d'\omterm{def:b3:rings:ideal}{idéal}
(\cref{prop:b3:spectral:ideal})~: si $S$ est \omterm{def:b3:spectral:compact}{compact} et $A,
B$ bornés, alors $ASB$ est \omterm{def:b3:spectral:compact}{compact}. En déduire que si $ST =
TS = I$ pour un $S$ borné, et $\dim H = \infty$, alors $T$
n'est pas \omterm{def:b3:spectral:compact}{compact} --- et réconcilier avec
\cref{exo:b3:spectral:1}(b).
\end{exercise}

\begin{exercise}[$\star\star$]\label{exo:b3:spectral:4}
(Hilbert--Schmidt) Soit $k \in L^2(\intcc01^2)$ et $(e_n)$ une
\omterm{def:b3:hilbert:onb}{base hilbertienne} de $L^2(\intcc01)$.
(a) Montrer que $\vertiii{T_k} \leq \norm k_{L^2}$
\emph{(Cauchy--Schwarz en la variable $y$, puis Tonelli)}.
(b) Développer $k(x,y) = \sum_{m,n}c_{mn}e_m(x)\overline{e_n(y)}$
dans $L^2$ du carré (justifier que les produits forment une
\omterm{def:b3:hilbert:onb}{base hilbertienne} là), et montrer que tronquer la somme donne
des opérateurs de rang fini convergeant vers $T_k$ en \omterm{def:b3:banach:operator}{norme
d'opérateur}~: $T_k$ est \omterm{def:b3:spectral:compact}{compact}.
\end{exercise}

\begin{exercise}[$\star\star$]\label{exo:b3:spectral:5}
Soit $T$ \omterm{def:b3:spectral:selfadjoint}{autoadjoint} avec $\langle x, Tx\rangle \geq 0$ pour
tout $x$ (opérateur \emph{positif}).
(a) Montrer que les valeurs propres sont $\geq 0$ et que
$\vertiii T = \sup_{\norm x\leq1}\langle x, Tx\rangle$.
(b) Prouver l'inégalité de Cauchy--Schwarz généralisée
$\abs{\langle x, Ty\rangle}^2 \leq \langle x, Tx\rangle\langle
y, Ty\rangle$.
\end{exercise}

\begin{exercise}[$\star\star$]\label{exo:b3:spectral:6}
Sur $L^2(\intcc01)$, soit $(Mf)(x) = x\,f(x)$.
(a) Montrer que $M$ est borné, \omterm{def:b3:spectral:selfadjoint}{autoadjoint}, avec $\vertiii M =
1$, mais n'a \emph{aucune} valeur propre.
(b) Montrer que $M$ n'est pas \omterm{def:b3:spectral:compact}{compact} \emph{(exhiber une suite
bornée dont l'image n'a aucune sous-suite convergente, p.\
ex.\ des indicateurs normalisés d'intervalles se rétrécissant
près de $1$ --- ou invoquer le théorème spectral)}.
(c) Où la preuve du~\cref{lem:b3:spectral:existence} casse-t-elle
pour $M$~?
\end{exercise}

\begin{exercise}[$\star\star$]\label{exo:b3:spectral:7}
(Volterra) Sur $L^2(\intcc01)$, soit $Vf(x) = \int_0^xf(y)\dd
y$.
(a) Montrer que $V$ est \omterm{def:b3:spectral:compact}{compact} (Hilbert--Schmidt à noyau
$\mathbf 1_{y < x}$) mais non \omterm{def:b3:spectral:selfadjoint}{autoadjoint}~; calculer $V^*$.
(b) Montrer que $V$ n'a aucune valeur propre non nulle.
\emph{(De $Vf = \lambda f$~: $f$ a un représentant continu,
puis est $\mathcal C^1$, et résout $\lambda f' = f$, $f(0) =
0$.)}
(c) Conclure que la compacité seule ne produit aucun vecteur
propre~: l'autoadjonction dans
\cref{thm:b3:spectral:spectral} est essentielle.
\end{exercise}

\begin{exercise}[$\star\star\star$]\label{exo:b3:spectral:8}
(Courant--Fischer) Soit $T$ \omterm{def:b3:spectral:compact}{compact}, \omterm{def:b3:spectral:selfadjoint}{autoadjoint},
\emph{positif}, de valeurs propres $\mu_1 \geq \mu_2 \geq
\dots > 0$ (répétées par multiplicités, vecteurs propres
$e_1, e_2, \dots$). Montrer~:
\[
\mu_{k} = \max_{\substack{V \subseteq H \\ \dim V = k}}\
\min_{\substack{x \in V\\ \norm x = 1}}\ \langle x, Tx\rangle
= \min_{\substack{W \subseteq H\\ \operatorname{codim}W =
k-1}}\ \max_{\substack{x\in W\\ \norm x = 1}}\ \langle x,
Tx\rangle .
\]
\emph{(Tester $V = \operatorname{Vect}(e_1,\dots,e_k)$~; pour
la borne supérieure intersecter tout $V$ avec des espaces de
type $\operatorname{Vect}(e_k, e_{k+1}, \dots)$~: le
comptage de dimension force une intersection non triviale.)}
En déduire que les valeurs propres dépendent monotonement de
$T$ ($T \leq S \Rightarrow \mu_k(T) \leq \mu_k(S)$).
\end{exercise}

\begin{exercise}[$\star\star$]\label{exo:b3:spectral:9}
En utilisant le~\cref{thm:b3:spectral:fredholm} pour $Tf(x) =
\int_0^1\min(x,y)f(y)\dd y$ (\cref{ex:b3:spectral:minxy})~:
pour quels $\lambda \in \R$ l'équation intégrale
\[
f(x) - \lambda\int_0^1\min(x,y)\,f(y)\,\dd y = g(x)
\]
a-t-elle une solution unique $f \in L^2$ pour tout $g \in
L^2$~? Que se passe-t-il aux valeurs exceptionnelles~?
\end{exercise}

\begin{exercise}[$\star\star\star$]\label{exo:b3:spectral:10}
Soit $S$ le décalage sur $\ell^2$
(\cref{exo:b3:banach:1}).
(a) Montrer que $S$ n'a aucune valeur propre, tandis que tout
$\lambda$ avec $\abs\lambda < 1$ est valeur propre de $S^*$
(trouver les vecteurs propres explicitement, suites
géométriques).
(b) Ni $S$ ni $S^*$ n'est \omterm{def:b3:spectral:compact}{compact}~: vérifier via un test à la
\cref{exo:b3:spectral:2} sur $(e_n)$.
(c) Commenter~: pour des opérateurs non \omterm{def:b3:spectral:selfadjoint}{autoadjoints}, non
\omterm{def:b3:spectral:compact}{compacts}, le paysage des valeurs propres peut être n'importe
quoi depuis le vide jusqu'à un disque plein --- la notion qui
survit est le \emph{spectre}, étudié dans un cours ultérieur.
\end{exercise}

\begin{exercise}[$\star\star$]\label{exo:b3:spectral:11}
(Racines carrées) Soit $T$ \omterm{def:b3:spectral:compact}{compact}, \omterm{def:b3:spectral:selfadjoint}{autoadjoint}, positif
($\langle x, Tx\rangle \geq 0$) sur un \omterm{def:b3:hilbert:inner}{espace de Hilbert} $H$,
avec décomposition spectrale $Tx = \sum_n\mu_n\langle e_n,
x\rangle e_n$ ($\mu_n > 0$).
(a) Définir $Sx = \sum_n\sqrt{\mu_n}\,\langle e_n, x\rangle
e_n$~; montrer que $S$ est \omterm{def:b3:spectral:compact}{compact}, \omterm{def:b3:spectral:selfadjoint}{autoadjoint}, positif, avec
$S^2 = T$.
(b) Prouver l'unicité~: tout $R$ \omterm{def:b3:spectral:compact}{compact} positif \omterm{def:b3:spectral:selfadjoint}{autoadjoint}
avec $R^2 = T$ préserve les espaces propres de $T$
\emph{($RT = R^3 = TR$~: $R$ commute avec $T$, donc
$R(\ker(T - \mu)) \subseteq \ker(T - \mu)$)}, et sur
$\ker(T - \mu)$, $R$ est un opérateur positif dont le carré
est $\mu\,\mathrm{id}$ sur un espace de dimension finie~: le
diagonaliser là et conclure $R = \sqrt\mu\,\mathrm{id}$ sur
chaque espace propre, d'où $R = S$.
(c) Calculer $\sqrt G$ pour l'opérateur de la corde $G$ du~\cref{pb:b3:spectral:1}~: quel noyau a les valeurs propres
$\frac1{n\pi}$ sur la \omterm{def:b3:topology:basis}{base} en sinus~? (Exprimer $\sqrt G$
comme limite $L^2$ de noyaux~; aucune forme fermée n'est
requise.)
\end{exercise}

\begin{exercise}[$\star\star\star$]\label{exo:b3:spectral:12}
(Décomposition en valeurs singulières) Soit $T \in \mathcal
L(H)$ \emph{\omterm{def:b3:spectral:compact}{compact}}, pas nécessairement \omterm{def:b3:spectral:selfadjoint}{autoadjoint}.
(a) Montrer que $T^*T$ est \omterm{def:b3:spectral:compact}{compact}, \omterm{def:b3:spectral:selfadjoint}{autoadjoint}, positif~;
soit $(e_n)$ une famille orthonormée de vecteurs propres avec
$T^*Te_n = s_n^2e_n$, $s_n > 0$ (les \emph{valeurs
singulières}), complétée par $\ker(T^*T) = \ker T$ (prouver
cette égalité).
(b) Poser $f_n = \frac{Te_n}{s_n}$~; montrer que $(f_n)$ est
orthonormée, et établir la \emph{SVD}~:
\[
Tx = \sum_n s_n\,\langle e_n, x\rangle\,f_n
\qquad (x \in H),
\]
avec convergence dans $H$.
(c) En déduire~: $\vertiii T = \max_ns_n$~; $T$ est limite en
\omterm{def:b3:banach:operator}{norme d'opérateurs} de rang fini (redémontrant le
\cref{prop:b3:spectral:ideal} en sens inverse pour les
\omterm{def:b3:hilbert:inner}{espaces de Hilbert})~; et pour l'opérateur de Volterra $V$ du~\cref{exo:b3:spectral:7}, qui n'a aucune valeur propre,
expliquer pourquoi la SVD existe néanmoins et quels en sont
les ingrédients (identifier $V^*V$ comme un opérateur à noyau
de type corde --- calculer ses valeurs propres explicitement
est le territoire du~\cref{exo:b3:spectral:9}).
\end{exercise}

\section{Problème~: la corde vibrante et
\texorpdfstring{$\zeta(2)$}{zeta(2)}}

\begin{problem}[{Problème du week-end --- opérateur de Green,
base en sinus, et formule de trace}]\label{pb:b3:spectral:1}
Nous résolvons le problème aux valeurs propres de la corde
vibrante à extrémités fixes --- $-u'' = \nu u$, $u(0) = u(1)
= 0$ --- par théorie des opérateurs, obtenons la \omterm{def:b3:topology:basis}{base}
orthonormée en sinus sans aucun calcul de Fourier, et
évaluons $\zeta(2)$ en comparant deux expressions pour la
\emph{trace} de l'opérateur de Green. On définit, sur
$L^2(\intcc01)$,
\[
(Gf)(x) = \int_0^1 g(x,y)\,f(y)\,\dd y,
\qquad
g(x, y) = \min(x,y)\,\bigl(1 - \max(x,y)\bigr).
\]

\medskip
\noindent\textbf{Partie I --- L'opérateur de Green.}
\begin{enumerate}
  \item Montrer que $g$ est \omterm{def:b3:topology:continuity}{continue}, symétrique, avec $0
        \leq g \leq \frac14$, et que $G$ est \omterm{def:b3:spectral:compact}{compact} et
        \omterm{def:b3:spectral:selfadjoint}{autoadjoint} (\cref{ex:b3:spectral:examples}(c)).
  \item Pour $f$ \omterm{def:b3:topology:continuity}{continue}, montrer que $u = Gf$ est $\mathcal
        C^2$ avec
        \[
        -u'' = f, \qquad u(0) = u(1) = 0
        \]
        \emph{(écrire $u(x) = (1-x)\int_0^xyf(y)\dd y +
        x\int_x^1(1-y)f(y)\dd y$ et dériver deux fois)}.
        Réciproquement, si $u \in \mathcal C^2$ avec $u(0) =
        u(1) = 0$, alors $G(-u'') = u$~: $G$ inverse
        l'opérateur de la corde.
  \item Montrer $\ker G = \{0\}$ \emph{(si $Gf = 0$ avec $f
        \in L^2$~: tester contre des $\varphi$ \omterm{def:b3:topology:continuity}{continues},
        transférer $G$ par symétrie/Fubini sur $\varphi$, et
        utiliser le lemme fondamental
        \cref{cor:b3:lp:fundlemma} --- ou régulariser)}, et
        que $G$ est un opérateur positif~: $\langle f,
        Gf\rangle \geq 0$. \emph{(Pour $f$ \omterm{def:b3:topology:continuity}{continue}~:
        $\langle f, Gf\rangle = \int_0^1 (u')^2$ avec $u =
        Gf$, par parties~; conclure par \omterm{ex:b3:lebesgue:gamma}{densité}.)}
\end{enumerate}

\medskip
\noindent\textbf{Partie II --- Diagonalisation~: la \omterm{def:b3:topology:basis}{base} en
sinus.}
\begin{enumerate}[resume]
  \item Montrer que les fonctions propres de $G$ à valeur
        propre $\lambda \ne 0$ sont, à scalaires près, les
        solutions de $-\lambda u'' = u$, $u(0) = u(1) = 0$
        \emph{(une fonction propre a un représentant continu
        --- $Gf$ est \omterm{def:b3:topology:continuity}{continue} pour $f \in L^2$, pourquoi~?
        --- donc est $\mathcal C^2$ par bootstrap de la
        question 2)}.
  \item Résoudre le problème aux limites~: les valeurs
        propres de $G$ sont $\lambda_n = \frac1{n^2\pi^2}$
        ($n \geq 1$), avec fonctions propres normalisées
        $e_n(x) = \sqrt2\,\sin(n\pi x)$~; vérifier
        l'orthonormalité par intégration directe comme test
        de cohérence.
  \item Conclure du~\cref{thm:b3:spectral:spectral} et de la
        question 3 que
        $\bigl(\sqrt2\sin(n\pi x)\bigr)_{n\geq1}$ est une
        \emph{\omterm{def:b3:topology:basis}{base}} orthonormée de $L^2(\intcc01)$ --- ni
        Stone--Weierstrass, ni série de Fourier nécessaires.
        Développer $f(x) = x(1-x)$ dans cette \omterm{def:b3:topology:basis}{base} et écrire
        \omterm{thm:b3:hilbert:parseval}{Parseval} pour elle.
\end{enumerate}

\medskip
\noindent\textbf{Partie III --- La formule de trace et
$\zeta(2)$.}
\begin{enumerate}[resume]
  \item Prouver les deux identités
        \[
        \langle e_n, Ge_n\rangle = \lambda_n
        \quad\text{et}\quad
        \sum_{n\geq1}\lambda_n = \int_0^1 g(x,x)\,\dd x .
        \]
        Pour la seconde (la \emph{formule de trace})~:
        développer $g(x, \cdot)$, pour $x$ fixé, dans la \omterm{def:b3:topology:basis}{base}
        $(e_n)$ --- montrer que les coefficients sont
        $\lambda_ne_n(x)$, de sorte que $g(x, \cdot) =
        \sum_n\lambda_ne_n(x)\,e_n$ dans $L^2$. Ici les sinus
        sont explicites~: vérifier directement que
        $\sum_n\lambda_ne_n(x)e_n(y)$ converge
        \emph{uniformément} sur le carré (comparer avec
        $\sum \frac2{n^2\pi^2}$), donc sa somme est \omterm{def:b3:topology:continuity}{continue}
        et, ayant les mêmes développements $L^2$ en $y$ pour
        chaque $x$, égale $g(x,y)$ partout. Poser $y = x$ et
        intégrer terme à terme.
  \item Calculer $\int_0^1g(x,x)\dd x = \int_0^1x(1-x)\dd x =
        \frac16$, et conclure
        \[
        \sum_{n\geq1}\frac{1}{n^2\pi^2} = \frac16,
        \qquad\text{i.e.}\qquad
        \boxed{\ \zeta(2) = \frac{\pi^2}{6}\ } :
        \]
        la somme d'Euler issue d'une trace d'opérateur.
  \item Redériver $\zeta(2)$ d'une troisième façon~: appliquer
        \omterm{thm:b3:hilbert:parseval}{Parseval} dans la \omterm{def:b3:topology:basis}{base} en sinus à la fonction
        constante $\mathbf 1$, calculer $\sum_{n \text{
        impair}}\frac1{n^2}$, et conclure. Puis comparer les
        mécanismes~: en quel sens l'argument de trace des
        questions 7--8 est-il «~\omterm{thm:b3:hilbert:parseval}{Parseval} appliqué au noyau
        entier d'un coup~»~?
\end{enumerate}

\medskip
\noindent\textbf{Partie IV --- La corde vibre.}
\begin{enumerate}[resume]
  \item (Séparation des variables, synthétisée) Pour $f \in
        L^2$, définir
        \[
        u(t, x) = \sum_{n\geq1}\;c_n\,
        \cos(n\pi t)\,\sqrt2\sin(n\pi x),
        \qquad c_n = \langle e_n, f\rangle .
        \]
        Montrer que la série converge dans $L^2(\intcc01)$
        pour chaque $t$, que $t\mapsto u(t, \cdot)$ est
        \omterm{def:b3:topology:continuity}{continue} à valeurs dans $L^2$, et que pour $f$ dans
        l'enveloppe d'un nombre fini de $e_n$ elle résout
        l'équation des ondes $\partial_t^2u =
        \partial_x^2u$ avec $u(0) = f$, $\partial_tu(0) =
        0$, extrémités fixes. Les valeurs propres $n^2\pi^2$
        sont les fréquences au carré~: les harmoniques de la
        corde --- expliquer l'interprétation musicale du~\cref{thm:b3:spectral:spectral} en un paragraphe.
\end{enumerate}

\medskip
\noindent\textbf{Partie V --- Dividendes variationnels~: la
méthode de la puissance, stabilité de Weyl, et une borne
rigoureuse sur $\pi$.} Soit $A$ un \omterm{def:b3:spectral:compact}{opérateur compact}
\omterm{def:b3:spectral:selfadjoint}{autoadjoint} \emph{positif} de valeurs propres $\mu_1 \geq
\mu_2 \geq \cdots > 0$ et vecteurs propres orthonormés
$(u_n)$~; $R(x) = \frac{\langle x, Ax\rangle}{\norm x^2}$.
Les formules min--max sont le~\cref{exo:b3:spectral:8}~; ici
on les dépense.
\begin{enumerate}[resume]
  \item (Méthode de la puissance) Pour $x \neq 0$ écrire $m_p
        = \sum_n\mu_n^p\abs{\langle u_n, x\rangle}^2$.
        Montrer $m_pm_{p+2} \geq m_{p+1}^2$
        (Cauchy--Schwarz), en déduire la chaîne
        \[
        R(x) \;\leq\; \frac{\langle Ax, Ax\rangle}{\langle x,
        Ax\rangle} \;\leq\; R(Ax) \;\leq\; \mu_1,
        \]
        et prouver que si $\langle u_1, x\rangle \neq 0$,
        alors $R(A^kx) \to \mu_1$~: itérer l'opérateur sur
        tout vecteur générique calcule la valeur propre
        maximale --- la méthode de la puissance de l'analyse
        numérique, certifiée.
  \item (Stabilité de Weyl) Pour $A, B$ \omterm{def:b3:spectral:compact}{compacts}
        \omterm{def:b3:spectral:selfadjoint}{autoadjoints} positifs, déduire du~\cref{exo:b3:spectral:8} que
        \[
        \abs{\mu_n(A) - \mu_n(B)} \;\leq\; \vertiii{A - B}
        \qquad\text{pour tout } n :
        \]
        le spectre entier est $1$-lipschitzien en la \omterm{def:b3:banach:operator}{norme
        d'opérateur} --- les valeurs propres de grands
        systèmes symétriques peuvent être calculées à partir
        d'approximations avec une erreur garantie.
  \item Appliquer la borne de Rayleigh à $G$ avec la fonction
        test $u(x) = x(1-x)$~: résoudre $-w'' = u$, $w(0) =
        w(1) = 0$ pour obtenir $Gu = w = \frac{x(1-x)(1 + x -
        x^2)}{12}$, calculer
        \[
        \norm u_2^2 = \frac1{30}, \qquad
        \langle u, Gu\rangle = \frac1{12}\Bigl(\frac1{30} +
        \frac1{140}\Bigr) = \frac{17}{5040},
        \qquad R(u) = \frac{17}{168},
        \]
        et conclure la borne \emph{rigoureuse} $\frac1{\pi^2}
        = \lambda_1 \geq \frac{17}{168}$, i.e.\ $\pi \leq
        \sqrt{168/17} < 3.1437$.
  \item Une étape de la chaîne de la question 11, sur la même
        fonction test~: en utilisant $\int_0^1(x - x^2)^4\dd
        x = \frac1{630}$, calculer
        \[
        \norm{Gu}_2^2 = \frac1{144}\Bigl(\frac1{30} +
        \frac2{140} + \frac1{630}\Bigr) = \frac{31}{90720},
        \qquad
        \frac{\langle Gu, Gu\rangle}{\langle u, Gu\rangle} =
        \frac{31}{306},
        \]
        et conclure $\pi \leq \sqrt{306/31} < 3.1419$~: deux
        intégrales, quatre chiffres corrects. (Chaque
        itération ultérieure élève approximativement au
        carré la précision~: l'écart spectral
        $\lambda_1/\lambda_2 = 4$ entraîne une convergence
        géométrique.)
\end{enumerate}

\medskip
\noindent\textbf{Partie VI --- La trace de $G^2$, et
$\zeta(4)$.}
\begin{enumerate}[resume]
  \item Montrer que $\sum_n\lambda_n^2 =
        \iint_{\intcc01^2}g(x,y)^2\,\dd x\,\dd y$
        \emph{(développer $g$ sur la \omterm{def:b3:topology:basis}{base} produit
        $(e_m(x)e_n(y))_{m,n}$ de $L^2(\intcc01^2)$ --- une
        \omterm{def:b3:hilbert:onb}{base hilbertienne}, cf.\ l'~\cref{exo:b3:spectral:5} ---
        et appliquer \omterm{thm:b3:hilbert:parseval}{Parseval} sur le carré~; la question 7
        identifie les coefficients)}.
  \item Calculer l'intégrale double~:
        \[
        \iint g^2 = 2\int_0^1(1-x)^2\Bigl(\int_0^x
        y^2\,\dd y\Bigr)\dd x
        = \frac23\int_0^1x^3(1-x)^2\,\dd x = \frac1{90} .
        \]
  \item Conclure $\zeta(4) = \dfrac{\pi^4}{90}$~; expliquer,
        sans calcul, comment les traces des puissances
        supérieures $G^k$ produisent $\zeta(2k) \in
        \pi^{2k}\,\Q$ pour tout $k \geq 1$, et pourquoi les
        valeurs impaires $\zeta(3), \zeta(5), \dots$ sont
        structurellement hors de portée de cette machine.
  \item ($\pi$ par en dessous) De $\lambda_1^2 \leq
        \sum_n\lambda_n^2 = \frac1{90}$ déduire $\pi \geq
        90^{1/4} > 3.080$, et assembler avec la question 14
        le verdict bilatéral
        \[
        3.080 \;<\; \pi \;<\; 3.1419,
        \]
        obtenu entièrement à partir de l'arithmétique de la
        corde vibrante. Quel côté converge plus vite si l'on
        utilise des traces plus élevées
        $(\operatorname{tr}G^{2k})^{-1/4k}$, et pourquoi~?
\end{enumerate}

\medskip
\noindent\textbf{Partie VII --- Forçage et résonance.} Fixer
$\nu \in \R$ et considérer la corde forcée $-u'' - \nu u =
f$, $u(0) = u(1) = 0$, avec $f \in L^2$ et $c_n = \langle
e_n, f\rangle$.
\begin{enumerate}[resume]
  \item Supposer $\nu \notin \{n^2\pi^2 : n \geq 1\}$. Montrer
        que
        \[
        u = \sum_{n\geq1}\frac{c_n}{n^2\pi^2 - \nu}\,e_n
        \]
        converge dans $L^2$ et uniformément sur $\intcc01$
        \emph{(Cauchy--Schwarz entre $(c_n)$ et les queues de
        $\sum n^{-4}$, avec $\norm{e_n}_\infty = \sqrt2$)},
        et qu'elle satisfait $u = Gf + \nu Gu$ --- la forme
        en coordonnées explicites de l'\omterm{thm:b3:spectral:fredholm}{alternative de
        Fredholm} (\cref{thm:b3:spectral:fredholm}), avec
        unicité.
  \item Supposer $\nu = m^2\pi^2$. Montrer que $u = Gf + \nu
        Gu$ a une solution ssi $c_m = 0$, unique à
        l'addition de multiples de $e_m$ près. Lecture
        physique~: pousser une balançoire exactement à sa
        propre fréquence.
  \item Pour $\nu < \pi^2$, montrer que l'opérateur de
        solution $R_\nu\colon f \mapsto u$ est borné sur
        $L^2$ de norme $\frac1{\pi^2 - \nu}$, \omterm{def:b3:spectral:compact}{compact},
        \omterm{def:b3:spectral:selfadjoint}{autoadjoint}, et positif~: toute l'analyse spectrale
        redémarre, décalée de $\nu$.
  \item (Synthèse) Compiler le dictionnaire de ce problème~:
        valeur propre $\leftrightarrow$ fréquence au carré
        (harmoniques)~; trace $\leftrightarrow$ $\zeta(2)$~;
        norme de Hilbert--Schmidt $\leftrightarrow$
        $\zeta(4)$~; \omterm{thm:b3:spectral:fredholm}{alternative de Fredholm}
        $\leftrightarrow$ résonance~; min--max
        $\leftrightarrow$ bornes variationnelles ($\pi <
        3.1437$ à partir d'un polynôme). Un seul opérateur
        intégral, cinq chapitres d'analyse encaissés.
\end{enumerate}

\medskip
\noindent\textbf{Partie VIII --- Trois derniers échos.}
\begin{enumerate}[resume]
  \item ($\zeta(6)$, pour rien) L'identité de \omterm{thm:b3:hilbert:parseval}{Parseval} de la
        question 6 a donné $\sum_{n\text{ impair}}n^{-6} =
        \frac{\pi^6}{960}$. Scinder $\zeta(6)$ en $n$
        impairs et pairs et conclure
        \[
        \zeta(6) = \frac{\pi^6}{945},
        \]
        sans nouvelle intégrale~: la machine de la question
        17 (traces de $G^3$) aurait produit la même valeur
        au prix d'un noyau itéré --- \omterm{thm:b3:hilbert:parseval}{Parseval} sur une
        fonction bien choisie est ici la voie moins chère.
  \item (L'état fondamental est positif) Soit $A$ un
        \omterm{def:b3:spectral:compact}{opérateur compact} \omterm{def:b3:spectral:selfadjoint}{autoadjoint} positif sur
        $L^2(\intcc01)$ donné par un noyau continu
        symétrique $k > 0$ sur $\intoo01^2$, de plus grande
        valeur propre $\mu_1$. Montrer~: (a) tout maximiseur
        du quotient de Rayleigh est une fonction propre de
        $\mu_1$~; (b) si $u$ en est une, alors
        $\langle\abs u, A\abs u\rangle \geq \langle u,
        Au\rangle$, avec inégalité \emph{stricte} si $u$
        prend les deux signes sur des ensembles de \omterm{def:b3:measure:measure}{mesure}
        positive --- donc $u$ a un signe constant p.p., et
        $u = \mu_1^{-1}Au$ ne s'annule jamais sur
        $\intoo01$~; (c) $\mu_1$ est une valeur propre
        \emph{\omterm{def:b3:groups:simple}{simple}}. Vérifier chaque affirmation sur $G$~:
        $e_1 = \sqrt2\sin(\pi x) > 0$, et chaque $e_n$, $n
        \geq 2$, étant orthogonal à $e_1$, doit changer de
        signe (il le fait~: $n - 1$ zéros intérieurs).
  \item (Distance au spectre, et le prix de la résonance)
        Pour $\nu \notin \{n^2\pi^2\}$, montrer que
        l'opérateur de solution $R_\nu$ de la question 19 est
        borné, \omterm{def:b3:spectral:selfadjoint}{autoadjoint}, \omterm{def:b3:spectral:compact}{compact}, avec
        \[
        \vertiii{R_\nu} =
        \frac1{\min_{n\geq1}\,\abs{n^2\pi^2 - \nu}}
        = \frac1{\operatorname{dist}\bigl(\nu,
        \{n^2\pi^2\}\bigr)},
        \]
        la norme étant atteinte sur le mode le plus proche.
        Puis quantifier la balançoire de la question 20~: un
        forçage avec $f = e_1$ en $\nu = (1 -
        \varepsilon)\pi^2$ produit $u =
        \frac{1}{\varepsilon\pi^2}\,e_1$, une amplification
        par $\frac1\varepsilon$ de la réponse statique $Ge_1
        = \frac1{\pi^2}e_1$ --- à un pour cent sous le
        fondamental ($\varepsilon = 10^{-2}$), la corde
        répond cent fois plus fort.
\end{enumerate}
\end{problem}
